\documentclass[12pt,a4paper]{article}
\usepackage[utf8]{inputenc}
\usepackage[T1]{fontenc}
\usepackage{amsmath,amssymb,amsthm}
\usepackage{physics}
\usepackage{geometry}
\usepackage{enumitem}
\usepackage{hyperref}
\usepackage{fancyhdr}
\usepackage{titlesec}

\geometry{margin=1in}
\pagestyle{fancy}
\fancyhf{}
\rhead{Lecture 2}
\lhead{Advanced Quantum Mechanics}
\cfoot{\thepage}

\titleformat{\section}{\large\bfseries}{}{0em}{}
\titleformat{\subsection}{\normalsize\bfseries}{}{0em}{}

\newtheorem{definition}{Definition}

\title{\textbf{Lecture 2: Symmetry, Degeneracy, and Angular Momentum} \\ \large Rotation Generators and the Angular Momentum Algebra}
\author{Based on Leonard Susskind's \textit{Advanced Quantum Mechanics}}
\date{}

\begin{document}
\maketitle
\thispagestyle{fancy}

% ============================================================
\section{1. Symmetry and Degeneracy}
% ============================================================

When a Hamiltonian possesses a symmetry $[U, H] = 0$, the energy eigenstates come in \emph{multiplets}---sets of states related by the symmetry operator that all share the same energy.

\begin{definition}[Degeneracy]
An energy level $E_n$ is \textbf{degenerate} if there exist multiple linearly independent eigenstates with the same eigenvalue $E_n$.  The number of such states is the \textbf{degree of degeneracy}.
\end{definition}

\noindent If $H\ket{E} = E\ket{E}$ and $[U, H] = 0$, then:
\[
    H\bigl(U\ket{E}\bigr) = U\bigl(H\ket{E}\bigr) = E\bigl(U\ket{E}\bigr).
\]
So $U\ket{E}$ is also an eigenstate with the same energy $E$.  Unless $U\ket{E}$ is proportional to $\ket{E}$, the level is degenerate.

\medskip
\noindent\textbf{Example:} For a rotationally invariant Hamiltonian, all states related by rotations share the same energy.  This leads to the characteristic $(2\ell+1)$-fold degeneracy of angular momentum multiplets.

% ============================================================
\section{2. Angular Momentum as a Generator of Rotations}
% ============================================================

\subsection{2.1 Classical angular momentum}

In classical mechanics, the angular momentum of a particle with position $\vb{r}$ and momentum $\vb{p}$ is:
\[
    \vb{L} = \vb{r} \times \vb{p}.
\]
In component form:
\[
    L_x = yp_z - zp_y, \qquad L_y = zp_x - xp_z, \qquad L_z = xp_y - yp_x.
\]
Angular momentum is the generator of rotations: a rotation by angle $\theta$ about axis $\hat{n}$ is generated by $\vb{L}\cdot\hat{n}$.

\subsection{2.2 Quantum angular momentum operators}

Promoting to operators via $p_i \to -i\hbar\,\partial/\partial x_i$:
\[
    \hat{L}_x = -i\hbar\left(y\frac{\partial}{\partial z} - z\frac{\partial}{\partial y}\right), \quad
    \hat{L}_y = -i\hbar\left(z\frac{\partial}{\partial x} - x\frac{\partial}{\partial z}\right), \quad
    \hat{L}_z = -i\hbar\left(x\frac{\partial}{\partial y} - y\frac{\partial}{\partial x}\right).
\]

A finite rotation by angle $\theta$ about the $z$-axis is:
\[
    R_z(\theta) = e^{-i\theta L_z/\hbar}.
\]

% ============================================================
\section{3. The Angular Momentum Commutation Relations}
% ============================================================

The fundamental commutation relations of angular momentum are:
\[
    \boxed{[L_x, L_y] = i\hbar L_z, \qquad [L_y, L_z] = i\hbar L_x, \qquad [L_z, L_x] = i\hbar L_y.}
\]
Compactly written using the Levi-Civita symbol:
\[
    [L_i, L_j] = i\hbar\,\epsilon_{ijk}\,L_k.
\]

These commutation relations encode the geometry of rotations: rotations about different axes do not commute, and the commutator of two rotation generators yields the third.

\subsection{3.1 The total angular momentum operator}

Define the operator:
\[
    L^2 = L_x^2 + L_y^2 + L_z^2.
\]
A crucial property is that $L^2$ commutes with every component:
\[
    [L^2, L_i] = 0 \quad \text{for } i = x, y, z.
\]
Therefore, we can simultaneously diagonalize $L^2$ and \emph{one} component of $\vb{L}$ (conventionally $L_z$).  However, the individual components do not commute with each other, so we cannot simultaneously know all three.

% ============================================================
\section{4. Raising and Lowering Operators}
% ============================================================

\begin{definition}[Ladder operators]
The angular momentum \textbf{raising} and \textbf{lowering} operators are:
\[
    L_+ = L_x + iL_y, \qquad L_- = L_x - iL_y.
\]
Note that $L_+^\dagger = L_-$.
\end{definition}

\subsection{4.1 Key commutation relations}

\begin{align}
    [L_z, L_+] &= +\hbar\, L_+, \\
    [L_z, L_-] &= -\hbar\, L_-, \\
    [L_+, L_-] &= 2\hbar\, L_z.
\end{align}

\subsection{4.2 Action on eigenstates}

If $L_z\ket{m} = m\hbar\ket{m}$, then:
\[
    L_z\bigl(L_+\ket{m}\bigr) = (m+1)\hbar\bigl(L_+\ket{m}\bigr), \qquad
    L_z\bigl(L_-\ket{m}\bigr) = (m-1)\hbar\bigl(L_-\ket{m}\bigr).
\]
So $L_+$ raises the $L_z$ eigenvalue by $\hbar$, and $L_-$ lowers it by $\hbar$.  This is the ``ladder'' that generates the full spectrum of $L_z$ eigenvalues within a multiplet.

% ============================================================
\section{5. The Spectrum of Angular Momentum}
% ============================================================

\subsection{5.1 Bounded spectrum}

Since $L^2 = L_z^2 + \tfrac{1}{2}(L_+L_- + L_-L_+)$ and all terms are non-negative in expectation, the eigenvalues of $L_z$ must be bounded:
\[
    -\ell \leq m \leq \ell,
\]
where $\ell$ is determined by $L^2\ket{\ell, m} = \hbar^2\ell(\ell+1)\ket{\ell, m}$.

\subsection{5.2 Integer and half-integer values}

The ladder terminates: there must exist a top state $\ket{\ell, \ell}$ with $L_+\ket{\ell, \ell} = 0$ and a bottom state $\ket{\ell, -\ell}$ with $L_-\ket{\ell, -\ell} = 0$.  The requirement that the ladder has an integer number of steps forces:
\[
    \ell = 0, \;\tfrac{1}{2}, \;1, \;\tfrac{3}{2}, \;2, \;\ldots
\]
For each $\ell$, the magnetic quantum number $m$ takes the $2\ell+1$ values:
\[
    m = -\ell, \;-\ell+1, \;\ldots, \;\ell-1, \;\ell.
\]

\subsection{5.3 Degeneracy from rotational symmetry}

For a rotationally invariant Hamiltonian ($[H, \vb{L}] = 0$), all $2\ell + 1$ states in a multiplet have the same energy.  The energy depends on $\ell$ but not on $m$.

\medskip
\noindent\textit{The algebraic structure developed here---commutation relations, ladder operators, bounded spectra---is one of the most powerful tools in quantum mechanics.  It will be applied to the hydrogen atom and the harmonic oscillator in the next lecture.}

\end{document}
